% ------------------------------------------------------------------------------
% Slide: Teoría de la Equivalencia Motora: Un Mismo Sonido, Múltiples Caminos
% ------------------------------------------------------------------------------
\begin{frame}{Teoría de la Equivalencia Motora: Un Mismo Sonido, Múltiples Caminos}
    \begin{columns}[c]
        \begin{column}{0.51\textwidth}
            \scriptsize
            \begin{block}{Principio Físico: Redundancia Mecánica}
                \begin{itemize}\setlength{\itemsep}{0pt}\setlength{\parskip}{0pt}
                    \item El cerebro busca un \textbf{objetivo acústico} ($F_1, F_2$), no activar un músculo fijo.
                    \item Con más variables mecánicas que frecuencias de salida, \textbf{múltiples combinaciones musculares producen el mismo sonido}.
                \end{itemize}
            \end{block}
            \vspace{0.03cm}
            \begin{block}{Cómo se Manifiesta en las Mediciones}
                \begin{itemize}\setlength{\itemsep}{0pt}\setlength{\parskip}{0pt}
                    \item \textbf{Vocal /i/:} El mismo $F_2$ alto se logra sonriendo con mejilla y menos lengua, o dejando la cara quieta y moviendo solo la lengua.
                    \item \textbf{Vocal /a/:} Existen dos familias acústicas: una \textbf{A vertical} (solo mandíbula) y una \textbf{A horizontal} (con tracción de mejillas).
                \end{itemize}
            \end{block}
            \vspace{0.03cm}
            \begin{alertblock}{Consecuencia para la Decodificación}
                \begin{itemize}\setlength{\itemsep}{0pt}\setlength{\parskip}{0pt}
                    \item Si el hablante cambia de estrategia motora, el sonido sale igual pero la señal superficial viaja a otra zona del espacio PCA.
                    \item \textbf{Dilema clave:} O se \textbf{fija el gesto externo} entrenando al usuario, o se \textbf{mide la lengua} (grado de libertad oculto).
                \end{itemize}
            \end{alertblock}
        \end{column}
        \begin{column}{0.47\textwidth}
            \centering
            \placeholderfigura{0.70\textheight}{Espacio 3D / PCA mostrando el desplazamiento por Equivalencia Motora: dos familias de A o cambio de estrategia en I entre sesiones}
        \end{column}
    \end{columns}
\end{frame}
